%This is my super simple Real Analysis Homework template

\documentclass{article}
\usepackage[utf8]{inputenc}
\usepackage[english]{babel}
\usepackage[]{amsthm} %lets us use \begin{proof}
\usepackage{amsmath}
\usepackage[]{amssymb} %gives us the character \varnothing

\title{Homework 6}
\author{Aine Zhang}
\date\today
%This information doesn't actually show up on your document unless you use the maketitle command below

\begin{document}
\maketitle %This command prints the title based on information entered above

%Section and subsection automatically number unless you put the asterisk next to them.
\section*{Problem 1}
Given function $f (x, y) = x + y$, find the maximum and minimum values of f over the domain
$${(x, y): 0 \leq y \leq x^2 \text{ and } - 1 \leq x \leq 1}$$

\begin{proof}
$$f(x,y)= x+y$$
$$g(x,y)=$$
Since $\nabla f = (1,1)$ is never zero, there are no interior critical points. So any extrema must be on the boundaries. \\
The boundary has three parts:
$y=0$, $-1 \le x \le 1$.
On this boundary, $f(x,0)=x$. Over $[-1,1]$, the minimum is $-1$ at $x=-1$, and the maximum is $1$ at $x=1$.
Upper boundary: $y=x^2$, $-1 \le x \le 1$.
Substitute into $f$: $f(x,x^2)=x+x^2.$ \\
Let $h(x)=x+x^2.$ Then $h'(x)=1+2x.$
Setting $h'(x)=0$ gives $1+2x=0  \Rightarrow  x=-\frac12.$
Then
\[
h\!\left(-\frac12\right)
=
-\frac12+\frac14
=
-\frac14. \qquad
h(-1)=0,  h(1)=2.
\]

On this boundary, the minimum is $-\frac14$ at 
$\left(-\frac12,\frac14\right)$, 
and the maximum is $2$ at $(1,1)$. \\
Vertical boundaries: $x=-1$ and $x=1$.
If $x=-1$, then $0 \le y \le 1$, and $f(-1,y)=-1+y.$
This is minimized at $y=0$ giving $-1$, 
and maximized at $y=1$ giving $0$. \\
If $x=1$, then $0 \le y \le 1$, and $f(1,y)=1+y.$

This is minimized at $y=0$ giving $1$, 
and maximized at $y=1$ giving $2$.

So, we get the minimum -1 at $(-1,0)$ and maximum value 2  $(1,1)$.

\end{proof}

\clearpage %Gives us a page break before the next section. Optional.

\section*{Exercise 4.8.7., Problem 2}

Let $x_1, x_2,.., x_n$ be positive real numbers. Prove the arithmetic-geometric mean inequality,
$$(x_1x_2 . . . x_n)^{\frac{1}{n}} \leq \frac{(x_1 + x_2 + ... + x_n)}{n} .$$
\textit{Hint}: Consider the function $f(x_1, x_2, ... , x_n) = \frac{x_1+x_2+...+x_n}{n}$ subject to the constraint $x_1x_2... x_n = c$, a constant.

\begin{proof}
By the hint, we construct function $f(x_1, x_2, ... , x_n) = \frac{x_1+x_2+...+x_n}{n}$ subject to the constraint $x_1x_2... x_n = c$, a constant. So 
$$f(x_1, x_2, ... , x_n) = \frac{x_1+x_2+...+x_n}{n}$$
$$ g(x_1, x_2, ... , x_n)=(x_1x_2 ... x_n)=c \Rightarrow (x_1x_2 ... x_n)-c =0$$
$$\nabla f= [\frac{1}{n}, \frac{1}{n}, ....\frac{1}{n}],  \nabla g=[x_2x_3...x_n,x_1x_3...x_n, ...x_1x_2...x_{n-1}]$$
$$\nabla f=\lambda \nabla g=\begin{cases}
    \lambda x_2x_3....x_n = \frac{1}{n} \\
    \lambda x_1x_3....x_n = \frac{1}{n}\\
    ... = \frac{1}{n} \\
    \lambda x_1x_2...x_{n-1}=\frac{1}{n}
\end{cases}$$
$$\lambda \neq 0 $$ so we know $x_2x_3....x_n=x_1x_3...x_n=x_1x_2...x_{n-1}$. Since all $x_1...x_n$ are positive reals, we can cancel terms, and we get the equality holds only if $x_1=x_2=...=x_n$. So, when all $x$'s are equal to each other, this is a critical point. \\
Let's say $x_1=x_2...=x_n=x$, then $\frac{x_1+x_2+...x_n}{n}=\frac{x+x+...+x}{n}=\frac{x_n}{n}=x$, and $x_1x_2...x_n=c \Rightarrow x^n=c, x =c^\frac{1}{n}$. \\
This point is a minimum, since the sum of a series of numbers is minimized when all the numbers are equal to each other. So, the critical point is a minimum of $f$ under the constraint $g.$ Thus, we know $\frac{x_1+x_2+...x_n}{n} \geq c ^\frac{1}{n}=(x_1....x_n)^\frac{1}{n}.$ Rearranging the equation gives us $(x_1x_2 . . . x_n)^{\frac{1}{n}} \leq \frac{(x_1 + x_2 + ... + x_n)}{n} .$
\end{proof}

\clearpage

\section*{Exercise 4.8.12., Problem 3}
Consider the plane defined by $P (x, y, z) = Ax + By + Cz + D = 0$ and the ellipsoid defined by $E(x, y, z) = \frac{x^2}{a^2}+ \frac{y^2}{b^2} +\frac{z^2}{c^2}-1=0$
\subsection*{i. }
Find conditions on $A, B, C, D, a, b, c$ such that the plane and the ellipsoid do not intersect.
\begin{proof}
The plane intersects the ellipsoid if and only if there exists a point on the ellipsoid satisfying $Ax+By+Cz = -D.$ \\
So, we can model this in a Lagrangian multiplier problem if we consider $L(x,y,z)=Ax+By+Cz$ restricted to the ellipsoid. 


    $f(A)\cdot f(B) \leq 0 \Leftrightarrow D^2-(A^2a^2+B^2b^2+C^2c^2) \leq 0$


Using Lagrange multipliers, one finds that the maximum and minimum
values of $L$ on the ellipsoid are
$$g_1=\frac{\sqrt{A^2a^2+B^2b^2+C^2c^2}+D}{\sqrt{A^2+B^2+C^2}}$$
$$g_2=\frac{-\sqrt{A^2a^2+B^2b^2+C^2c^2}+D}{\sqrt{A^2+B^2+C^2}}$$
Therefore, $Ax+By+Cz$ takes all values between $-\sqrt{A^2a^2+B^2b^2+C^2c^2}$ and $\sqrt{A^2a^2+B^2b^2+C^2c^2}.$
The plane intersects the ellipsoid if and only if $-D$ lies in this interval, which is
\[
|D| \le \sqrt{A^2a^2+B^2b^2+C^2c^2}.
\]

The plane and ellipsoid do not intersect when
\[
|D| > \sqrt{A^2a^2+B^2b^2+C^2c^2}.
\]

\end{proof}
\subsection*{ii. }
Find the minimum distance between the plane and the ellipsoid when they do not intersect.

\begin{proof}

The distance from a point $(x,y,z)$ to the plane is$
d(x,y,z)=\frac{|Ax+By+Cz+D|}{\sqrt{A^2+B^2+C^2}}.$

Since we the plane and ellipsoid do not intersect, the sign of 
$Ax+By+Cz+D$ is constant, so we can minimize $f(x,y,z)=Ax+By+Cz+D$ subject to the constraint $
g(x,y,z)=\frac{x^2}{a^2}+\frac{y^2}{b^2}+\frac{z^2}{c^2}-1=0.
$ \\

Using Lagrange multipliers, we compute

$$
\nabla f = (A,B,C), \qquad \nabla g = \left(\frac{2x}{a^2},\frac{2y}{b^2},\frac{2z}{c^2}\right).
$$

We set $\nabla f=\lambda \nabla g,$ so

$$
(A,B,C)=\lambda\left(\frac{2x}{a^2},\frac{2y}{b^2},\frac{2z}{c^2}\right).
$$

This gives

$$
x=\frac{A a^2}{2\lambda}, 
y=\frac{B b^2}{2\lambda}, 
z=\frac{C c^2}{2\lambda}. \quad \Rightarrow \quad 
\frac{A^2 a^4}{4\lambda^2 a^2}
+
\frac{B^2 b^4}{4\lambda^2 b^2}
+
\frac{C^2 c^4}{4\lambda^2 c^2}
=1.
$$

Simplifying,

$$
\frac{A^2 a^2 + B^2 b^2 + C^2 c^2}{4\lambda^2}=1.
$$

So, $4\lambda^2=A^2 a^2 + B^2 b^2 + C^2 c^2,$ and

$$
2\lambda=\pm \sqrt{A^2 a^2 + B^2 b^2 + C^2 c^2}.
$$

Now compute the extreme values of $Ax+By+Cz$:

$$
Ax+By+Cz
=
\frac{A^2 a^2 + B^2 b^2 + C^2 c^2}{2\lambda}
=
\pm \sqrt{A^2 a^2 + B^2 b^2 + C^2 c^2}.
$$

Therefore the minimum value of $|Ax+By+Cz+D|$ on the ellipsoid is

$$
|D|-\sqrt{A^2 a^2 + B^2 b^2 + C^2 c^2},
$$

since that's also what we assumed in part (i). \\
Hence the minimum distance between the plane and the ellipsoid is
$$
\frac{|D|-\sqrt{A^2 a^2 + B^2 b^2 + C^2 c^2}}{\sqrt{A^2+B^2+C^2}}.
$$
\end{proof}
\clearpage

\section*{Exercise 4.8.14}
 Consider two nonparallel planes in $\mathbb{R}^3$. Find the point on their line of intersection closest to the origin in $\mathbb{R}^3$.

\begin{proof}
     Lets say there are two nonparallel planes $M$ and $N$, where $$M(x,y,z)=Ax+By+Cz=D, N(x,y,z)=Ex+Fy+Gz=H.$$ The line where they intersect is when both equations are satisfied and share the same $x,y,z$. We want to find the smallest distance between $P$ and $(0,0,0).$\\
     Let $f$ be the distance formula $$f(x,y,z)=x^2+y^2+z^2,$$ and $M$ and $N$ act as constraints on $f.$\\
     We want to use Lagrangian multipliers, so we find $$\nabla f=[2x, 2y, 2z], \nabla M =[A, B, C], \nabla N=[E,F,G]$$
     By applying the multipliers, we want
     $ \nabla f = \lambda_1 \nabla M + \lambda_2\nabla N$, which leads us to 
     $$\begin{bmatrix}
         2x \\
         2y \\
         2z \\
 \end{bmatrix} =\lambda_1 \begin{bmatrix}
         A \\
         B \\
         C \\
     \end{bmatrix}+\lambda_2\begin{bmatrix}
         E \\
         F\\
         G
     \end{bmatrix}$$
     In order to solve for $\begin{bmatrix}
         x & y &z
     \end{bmatrix},$ we want to plug in $M$ and $N$ into $\begin{bmatrix}
         x & y & z
     \end{bmatrix}=\lambda_1 \begin{bmatrix}
         A & B & C
     \end{bmatrix} + \lambda_2 \begin{bmatrix}
         E & F & G
     \end{bmatrix}. $
     If we denote $$a = [A,B,C] \cdot [A, B, C], b =[A,B,C]\cdot [E,F,G], c=[E,F,G]\cdot [E,F,G], $$
     plugging this into the equations we obtain $a\lambda_1+b\lambda_2=D, b\lambda_1+c\lambda_2=H\Rightarrow \lambda_1=\frac{cD-bH}{ac-b^2}, \lambda_2=\frac{-bD+aH}{ac-b^2}$
     Since the planes are nonparallel, their normal vectors $(A,B,C)$ and $(E,F,G)$ are not multiples of one another, so $ac-b^2 \neq 0$.
     So the point on the line of intersection closest to the origin where
     $$
     P
     =
     \lambda_1(A,B,C)
     +
     \lambda_2(E,F,G),
     $$
     where
     $$
     \lambda_1=\frac{cD-bH}{ac-b^2}, 
     \qquad 
     \lambda_2=\frac{-bD+aH}{ac-b^2}.
     $$
     Solving this gives us the coordinates 
    
    \[
    x
    =
    \frac{
    (cD - bH)A + (aH - bD)E
    }{ac - b^2},
    \]
    
    \[
    y
    =
    \frac{
    (cD - bH)B + (aH - bD)F
    }{ac - b^2},
    \]
    
    \[
    z
    =
    \frac{
    (cD - bH)C + (aH - bD)G
    }{ac - b^2}.
    \]
\end{proof} 


\end{document}